% secnumdepth=2：编号到 subsection，即模板中的三级标题
% tocdepth=2：目录显示到 subsection，即目录中显示三级标题
% 后三句：强制编号格式改成阿拉伯数字层级形式
\setcounter{secnumdepth}{2}
\setcounter{tocdepth}{2}

\renewcommand{\thechapter}{\arabic{chapter}}
\renewcommand{\thesection}{\thechapter.\arabic{section}}
\renewcommand{\thesubsection}{\thesection.\arabic{subsection}}

\ctexset{
  chapter = {
    name = {,},
    number = \arabic{chapter}
  }
}

% 使用 ctexset 对章节标题、页眉页脚等样式进行定制，以满足学校对字体和字号的要求。
% 一级标题（章标题）使用黑体、小三号字号，并居中排版；二级标题使用黑体四号；
% 三级标题使用黑体小四号；四级标题使用楷体小四号；五级标题使用宋体加粗小四号。
\ctexset{
  chapter={
    format=\centering\heiti\zihao{-3},
    nameformat=\heiti\zihao{-3},
    beforeskip=0pt,
    afterskip=12pt,
  },
  section={
    format=\raggedright\setstretch{1.5}\heiti\zihao{4},
    beforeskip=0pt,
    afterskip=0pt,
  },
  subsection={
    format=\raggedright\setstretch{1.5}\heiti\zihao{-4},
    beforeskip=0pt,
    afterskip=0pt,
  },
  subsubsection={
    format=\raggedright\setstretch{1.5}\kaishu\zihao{-4},
    beforeskip=0pt,
    afterskip=0pt,
  }
}

% 模板中的四级标题：编号为（1）、（2）、（3）……，楷体小四号，顶格。
\newcounter{fourthsection}[subsection]
\renewcommand{\thefourthsection}{（\arabic{fourthsection}）}
\newcommand{\fourthsection}[1]{%
  \refstepcounter{fourthsection}%
  \par\noindent{\setstretch{1.5}\kaishu\zihao{-4}\thefourthsection#1\par}%
  \nobreak
}

% 页面设置
\geometry{
  a4paper,
  top=3cm,
  bottom=4cm,
  left=2.5cm,
  right=2.5cm
}

% 页眉页脚位置
\setlength{\headheight}{14pt}   % 页眉区域高度，给 fancyhdr 留够空间
\setlength{\headsep}{2.01cm}    % 使页眉基线大致落在距顶边 1.5cm 附近
\setlength{\footskip}{1.0cm}    % 使页脚基线大致落在距底边 1.5cm 附近

% 页眉和页脚
\pagestyle{fancy}
\fancyhf{}

% 页眉：宋体五号（10.5pt）
\fancyhead[C]{\songti\zihao{5} 成都大学本科毕业设计（论文）}

% 页脚：宋体小五（9pt）
\fancyfoot[C]{\songti\zihao{-5}\thepage}

\renewcommand{\headrulewidth}{0.4pt}
\renewcommand{\footrulewidth}{0pt}

% 章节首页也使用相同页眉页脚
\makeatletter
\def\ps@plain{\ps@fancy}
\makeatother

% 正文字体及行距设置
\newcommand{\ThesisBodyFont}{%
  \songti\zihao{-4}%
  \linespread{1.25}\selectfont%
}

\newcommand{\ThesisBodyStyle}{%
  \ThesisBodyFont
  \setlength{\parindent}{2em}%
  \setlength{\parskip}{0pt}%
}

\makeatletter
\renewcommand{\normalsize}{\ThesisBodyFont}
\makeatother

\AtBeginDocument{%
  \ThesisBodyStyle
}

% 调研列表：编号格式为“（1）”“（2）”。
\newlist{researchlist}{enumerate}{1}
\setlist[researchlist]{%
  label=（\arabic*）,
  ref=（\arabic*）,
  labelindent=2em,
  leftmargin=*,
  labelsep=0.5em,
  itemsep=0pt,
  topsep=0pt,
  parsep=0pt,
  partopsep=0pt
}

\newcommand{\researchitem}[2]{%
  \item #1%
  \par\begingroup
    \leftskip=-\leftmargin
    \rightskip=0pt
    \parindent=2em
    \noindent\hspace*{2em}\ignorespaces#2\unskip\par
  \endgroup
}

% 目录点线：数字越小，点越密。默认约为 4.5，这里调成接近 Word 模板的密度。
\makeatletter
\renewcommand{\@dotsep}{1}
\renewcommand*\l@chapter[2]{%
  \ifnum \c@tocdepth >\m@ne
    \addpenalty{-\@highpenalty}%
    \vskip 1.0em \@plus\p@
    {\bfseries\@dottedtocline{0}{0em}{2.0em}{#1}{#2}}%
    \penalty\@highpenalty
  \fi
}
\makeatother

% 图题、表题格式：居中，黑体，五号。
\makeatletter
\@ifpackageloaded{caption}{%
  \DeclareCaptionFont{thesiscaptionfont}{\heiti\zihao{5}}
  \captionsetup{
    font=thesiscaptionfont,
    labelfont=thesiscaptionfont,
    textfont=thesiscaptionfont,
    justification=centering,
    singlelinecheck=true,
    parindent=0pt
  }
  \captionsetup[figure]{position=bottom}
  \captionsetup[table]{position=top}
}{}
\makeatother

% 表内文字：中文宋体五号，数字和英文使用 Times New Roman。
\AtBeginEnvironment{tabular}{\songti\zihao{5}\fontspec{Times New Roman}}

% 算法块输入、输出标签。
\makeatletter
\@ifpackageloaded{algorithm}{\floatname{algorithm}{算法}}{}
\@ifpackageloaded{algpseudocode}{%
  \algrenewcommand\algorithmicrequire{\textbf{输入：}}
  \algrenewcommand\algorithmicensure{\textbf{输出：}}
}{}
\makeatother

% 浮动体内部不继承正文首行缩进，避免图片、表格和题名整体超出版心。
\AtBeginEnvironment{figure}{\setlength{\parindent}{0pt}}
\AtBeginEnvironment{table}{\setlength{\parindent}{0pt}}
\AtBeginEnvironment{algorithm}{\setlength{\parindent}{0pt}}

% 参考文献正文：宋体五号，1.25 倍行距。
\AtBeginEnvironment{thebibliography}{%
  \songti\zihao{5}%
  \linespread{1.25}\selectfont%
}

\newcommand{\makeAbstractPersonalInfo}[4]{%
  {%
    \setstretch{1.25}\kaishu\zihao{-4}
    \noindent\hspace*{\abstractPersonalInfoIndent}% 摘要处个人信息左边距
    \begin{tabular}{@{}l@{\hspace{3em}}l@{}}
      专业：#1 & 学号：\TNR{#2} \\[1.25em]
      学生：#3 & 指导教师：#4 \\[1.25em]
    \end{tabular}
  }\\
}

\newcommand{\makeChineseAbstractHeader}[5]{%
  \clearpage
  \begin{center}
    {\setstretch{1.5}\heiti\zihao{-2}#1\par}
  \end{center}
  \vspace{12pt}
  %\addcontentsline{toc}{chapter}{摘要}
  \makeAbstractPersonalInfo{#2}{#3}{#4}{#5}
}

\newcommand{\ChineseAbstractTextStart}{%
  \noindent {\heiti\zihao{-4} 摘要}：\begingroup\songti\zihao{-4}\ignorespaces
}

\newcommand{\ChineseAbstractTextFinish}{%
  \endgroup
}

\newcommand{\makeChineseKeywords}[1]{%
  \par\vspace{\baselineskip}\noindent{\heiti\zihao{-4} 关键词：}\,{\songti\zihao{-4} #1}
}

\newcommand{\makeEnglishAbstractPersonalInfo}[4]{%
  {%
    \setstretch{1.25}\fontspec{Times New Roman}\zihao{-4}
    \noindent\hspace*{\englishAbstractPersonalInfoIndent}%
    \begin{tabular}{@{}l@{\hspace{3em}}l@{}}
      Major：#1 & Student ID：\TNR{#2} \\[1.25em]
      Student：#3 & Instructor：#4 \\
    \end{tabular}
  }\\
}

\newcommand{\makeEnglishAbstractHeader}[5]{%
  \clearpage
  \begin{center}
    {\setstretch{1.25}\fontspec{Times New Roman}\bfseries\zihao{-2}#1\par}
  \end{center}
  \vspace{12pt}
  \makeEnglishAbstractPersonalInfo{#2}{#3}{#4}{#5}
}

\newcommand{\EnglishAbstractTextStart}{%
  {\heiti\zihao{-4} }\par\vspace{1ex}
  \noindent {\fontspec{Times New Roman}\bfseries\zihao{-4} Abstract}:%
  \begingroup\fontspec{Times New Roman}\zihao{-4}\setlength{\parindent}{2em}\ignorespaces
}

\newcommand{\EnglishAbstractTextFinish}{%
  \endgroup
}

\newcommand{\makeEnglishKeywords}[1]{%
  \par\vspace{\baselineskip}\noindent{\fontspec{Times New Roman}\bfseries\zihao{-4} Key words:}\,{\fontspec{Times New Roman}\zihao{-4} #1}
}
